\section{微扰论中的 QCD 梯度流}\label{sec:pert}

下文我们在 $D$ 维欧氏时空中工作，$D=4-2\ep$。梯度流形式把常规\footnote{我们
  将用``流动''与``常规''QCD 这两个词来区分定义在 $t>0$ 与定义在 $t=0$ 的量。}
QCD 的胶子场与夸克场 $A^a_\mu(x)$、$\psi(x)$ 延拓为 $(D+1)$ 维场
$B^a_\mu(t,x)$ 与 $\chi(t,x)$，其边界条件为
\begin{equation}
  \begin{split}
    B_\mu^a (t=0,x) = A_\mu^a (x)\,,\qquad \chi (t=0,x)= \psi(x)\,,
    \label{eq:bound}
  \end{split}
\end{equation}
流方程为
\begin{equation}
  \begin{split}
    \partial_t B_\mu^a &= \mathcal{D}^{ab}_\nu G_{\nu\mu}^b + \kappa
    \mathcal{D}^{ab}_\mu \partial_\nu B_\nu^b\,,\\ \partial_t \chi &= \Delta \chi
    - \kappa \partial_\mu B_\mu^a T^a \chi\,,\\ \partial_t \overline
    \chi &= \overline \chi \overleftarrow \Delta + \kappa \overline \chi
    \partial_\mu B_\mu^a T^a\,,
    \label{eq:flow}
  \end{split}
\end{equation}
其中``流时间''$t$ 是质量量纲为 $-2$ 的参数，$\kappa$ 是一个额外的规范参数，
它从物理可观测量中消去。

$(D+1)$ 维场强张量定义为
\begin{align}
  G_{\mu\nu}^a = \partial_\mu B_\nu^a - \partial_\nu B_\mu^a + f^{abc}
  B_\mu^b B_\nu^c\,,
\end{align}
伴随表示中的协变导数为
\begin{align}
  \mathcal{D}_\mu^{ab} = \delta^{ab} \partial_\mu - f^{abc} B_\mu^c\,,
\end{align}
且
\begin{equation}
  \begin{split}
    \Delta = (\partial_\mu + B_\mu) (\partial_\mu + B_\mu)\,,\qquad
    \overleftarrow{\Delta} = (\overleftarrow\partial\!_\mu - B_\mu)
    (\overleftarrow\partial\!_\mu - B_\mu)\,.
  \end{split}
\end{equation}
照例，伴随表示的色指标记作 $a,b,c,\ldots$，而 $\mu,\nu,\rho,\ldots$
为 $D$ 维洛伦兹指标。除非出于清晰起见，基本表示的色指标在全文中省略。
对称性生成元 $T^a$ 满足对易关系
\begin{equation}
  \begin{split}
    [T^a,T^b] = f^{abc}T^c\,,
  \end{split}
\end{equation}
其中 $f^{abc}$ 为结构常数。

流场方程使规范场组态在有限流时间 $t>0$ 处被涂抹。因此 $t>0$ 处的复合算符
不需要超出拉格朗日量的参数与场之外的重正化。对强耦合与夸克质量，其重正常数
与 $t=0$ 时相同；如 \citere{Luscher:2011bx} 所指出，流动胶子场在有限流时间
无需重正化。流动夸克场到 \nnlo\ 的重正常数是本文的一个副产品，将在下文给出。
